\scanpage{360}
\mspage{EZ 3}
Les applications shuffles et A-W sont \uline{assoc}

\(\left(\begin{aligned}
&\text{et \uline{shuffles} est (anti)-commutative}\\
&\text{A-W \quad - \quad est seulement comm à homotopie près.}
\end{aligned}\right.\)

On peut exprimer ceci avec des diagrammes \unsure{d'objets} \unsure{trisimpliciaux}

\margin{Ref Illusie \uline{thèse.}}
\diagram{\begin{tikzcd}[ampersand replacement=\&]
X_{***}\arrow[r,"N_{\unsure{123|12|**}}"]\arrow[d,"\Delta_{12}"'] \& X_{\bullet\bullet\bullet}\arrow[r,"\int_{12}"] \& X_{\bullet\bullet}\arrow[d,"\int"]\\
X_{**}\arrow[r,"N"] \& X_{\bullet\bullet} \& X_\bullet
\end{tikzcd}}

Clair pour \(\underline{\mathbb ZX}\).
\[
\deleted{\unsure{\mathbb ZX\otimes\mathbb ZX\otimes\mathbb ZX\otimes\mathbb ZX}}
\]
L'enoncé de commutativité s'enonce comme suit : \hfill \(X\) ens simpl
\[
\begin{aligned}
&\deleted{\mathrm{Hom}(\mathbb ZX\otimes\mathbb ZX,\mathbb Z(X\times X))}\\
&\mathrm{Hom}^{\unsure{\bullet|'}}\bigl(\widetilde{\mathbb Z(X\times X)},\widetilde{\mathbb ZX}\otimes\widetilde{\mathbb ZX}\bigr)\longrightarrow\mathbb Z
\end{aligned}
\]
\(\Sigma_2\) agit sur ce complexe
\[
\begin{aligned}
&\nearrow\quad\mathrm{Hom}\bigl(\mathbb Z(X_0\times X_0)\xrightarrow{\sim}\mathbb ZX_0\otimes\mathbb ZX_0\bigr)\longrightarrow\mathbb Z\\
&\hspace{4em} X\longrightarrow\mathrm{pt}\ \deleted{\longrightarrow}\\
&\mathrm{Hom}\bigl(\mathbb Z(\ ,\ ),\widetilde{\mathbb Z(\ )}\otimes\widetilde{\mathbb Z(\ )}\bigr)\longrightarrow\mathbb Z
\end{aligned}
\]
\uline{Theorème} (Dold) : \textquotedblleft{} \unsure{Ueber|Über} Steenrodsche operationen \ldots\textquotedblright

Pour tout \(X\), Ce complexe
\[
\left.\begin{aligned}
\mathrm{Hom}^{\unsure{'|\bullet}}\bigl(\mathbb Z(X\ldots\times X),\mathbb Z(X)\otimes\ldots\otimes\mathbb ZX\\
\mathrm{Hom}^{\unsure{'|\bullet}}\bigl(\mathbb ZX\otimes\ldots\otimes\mathbb ZX,\mathbb Z(X\times\ldots\times X)
\end{aligned}\right\}
\]
est une \unsure{resolution|résolution} libre de \(\mathbb Z\) dans la cat des \(\Sigma_n\)-modules

\inserted{pour \(X\) variable}

On tronque pour avoir un complexe de chaines
% archive sheet number: 361
